\section{Introduction}

\subsection{Background of the Study}

The global transportation sector is undergoing a significant transformation following the increasing adoption of electric vehicles (EVs). This shift is primarily driven by the need to reduce greenhouse gas emissions and dependence on fossil fuels(in this case diesel).

However, the deployment of EV charging infrastructure presents major challenges, particularly in regions with weak or unstable electrical grids. Such grids are often characterized by limited generation capacity, frequent power outages, and reliance on fossil-fuel-based backup systems.

To address these challenges, solar-assisted EV charging systems have been proposed. These systems integrate photovoltaic (PV) generation and battery energy storage to reduce dependence on conventional grid power. Despite their advantages, these systems still face difficulties in efficiently managing energy flow between solar generation, storage, grid supply, and EV charging demand.

In many real-world scenarios, particularly in developing regions, backup diesel generators are still used to ensure charging reliability. This introduces additional operational costs and increases carbon emissions, undermining the environmental benefits of EV adoption.

\subsection{Problem Statement}

Weak electrical grids face significant challenges in supporting reliable EV charging infrastructure while minimizing dependence on fossil-fuel-based backup systems. This leads to increased carbon emissions, higher operational costs, and reduced system efficiency.

Existing energy management approaches often assume stable grid conditions and do not adequately address the combined challenges of renewable energy integration, grid instability, and emissions reduction.

\subsection{Research Questions}

This study addresses the following main question:

How can an adaptive energy management system improve the reliability, efficiency, and sustainability of solar-assisted EV charging under weak grid conditions?

\subsection{Objectives}

The main objective of this study is to design and evaluate an adaptive energy management strategy for solar-assisted EV charging systems under weak grid conditions.

The specific objectives are:

\begin{itemize}
    \item To develop a mathematical model of a hybrid solar-battery-grid EV charging system.
    \item To design an adaptive energy management algorithm for optimal power allocation.
    \item To minimize fossil fuel consumption in backup generation systems.
    \item To reduce carbon emissions associated with EV charging.
    \item To evaluate system performance under varying solar and grid conditions.
\end{itemize}
